% (c) Nikita Lisitsa, lisyarus@gmail.com, 2026

\documentclass[10pt]{beamer}

\usepackage[T2A]{fontenc}
\usepackage[russian]{babel}
\usepackage{minted}

\usepackage[most]{tcolorbox}
\tcbuselibrary{minted}

\usepackage{graphicx}
\graphicspath{ {./images/} }

\usepackage{adjustbox}

\usepackage{color}
\usepackage{soul}
\usepackage{multicol}
\usepackage{multirow}
\usepackage{hyperref}
\usepackage{tikz}
\usetikzlibrary{arrows.meta,positioning,calc}

\usetheme{metropolis}

\definecolor{red}{rgb}{1,0,0}
\definecolor{green}{rgb}{0,0.5,0}
\definecolor{codebg}{RGB}{29,35,49}
\setminted{fontsize=\footnotesize}

\makeatletter
\newcommand{\slideimage}[1]{
  \begin{figure}
    \begin{adjustbox}{width=\textwidth, totalheight=\textheight-2\baselineskip-2\baselineskip,keepaspectratio}
      \includegraphics{#1}
    \end{adjustbox}
  \end{figure}
}
\makeatother

\title{Язык программирования C}
\subtitle{Лекция 5: стандартная библиотка, ввод-вывод}
\author{Лисица Никита}
\date{2026}

\setbeamertemplate{footline}[frame number]

\usemintedstyle{tango}

\begin{document}

\frame{\titlepage}

\begin{frame}[fragile]
\frametitle{Стандартная библиотека}
\begin{itemize}
\item Набор полезной функциональности, поставляемый вместе с компилятором языка C:
\begin{itemize}
\item \mintinline{cpp}|stdio.h| -- ввод/вывод (файлы, клавиатура, экран)
\item \mintinline{cpp}|stdlib.h| -- выделение памяти, различные алгоритмы
\item \mintinline{cpp}|string.h| -- работа со строками, массивами
\item \mintinline{cpp}|math.h| -- математические функции
\item \mintinline{cpp}|time.h| -- работа со временем
\end{itemize}
\pause
\item Код стандартной библиотеки десятки лет тестировался и улучшался \(\Longrightarrow\) лучше использовать её функции, а не писать свои (разве что в образовательных целях)
\end{itemize}
\end{frame}

\begin{frame}[fragile]
\frametitle{Стандартная библиотека}
\begin{itemize}
\item Сам бинарный файл стандартной библиотеки называется по-разному в разных системах:
\begin{itemize}
\item Linux: \mintinline{text}|libc.so|
\item macOS: \mintinline{text}|libSystem.B.dylib|
\item Windows: \mintinline{text}|ucrtbase.dll|
\end{itemize}
\pause
\item Компилятор сам это знает, нам его подставлять руками не нужно :)
\end{itemize}
\end{frame}

\begin{frame}[fragile]
\frametitle{О работе с устройствами}
\begin{itemize}
\item Память, клавиатура, экран, и т.п. -- отдельные от CPU устройства
\pause
\item Операционная система не даёт к ним прямого доступа, вместо этого мы можем использовать \textit{системные вызовы} с помощью особой инструкции процессора (\mintinline{cpp}|syscall| на x86\_64)
\pause
\item Делать системные вызовы руками очень неудобно (для этого нужно писать на языке ассемблера), и стандартная библиотека оборачивает их в более удобный для программирования интерфейс
\pause
\item Например, чтение файла может выполняться так:
\pause
\begin{itemize}
\item Программа на C: \mintinline{cpp}|fread(...)|
\pause
\item \(\rightarrow\) Стандартная библиотека: \mintinline{cpp}|read syscall|
\pause
\item \(\rightarrow\) Операционная система: обращается напрямую к устройству
\end{itemize}
\end{itemize}
\end{frame}

\begin{frame}[fragile]
\frametitle{Работа с файлами}
\begin{itemize}
\item Стандартная библиотека C работает с файлами как с указателями на специальные структуры \mintinline{cpp}|FILE*|
\pause
\item \mintinline{cpp}|struct FILE| внутри хранит
\pause
\begin{itemize}
\item \textit{Файловый дескриптор} -- число, идентифицирующее открытый файл в ОС
\pause
\item Временный буфер, использующийся для чтения/записи (чтобы не читать/писать по одному байту)
\pause
\item Текущее состояние -- положение в файле, состояние ошибки, флаг конца файла
\end{itemize}
\pause
\item Напрямую доступа к этим полям нет, нужно использовать соответствующие функции из \mintinline{cpp}|stdio.h|
\end{itemize}
\end{frame}

\begin{frame}[fragile]
\frametitle{Работа с файлами}
\usemintedstyle{lightbulb}
\begin{tcblisting}{listing engine=minted, minted language=cpp, minted options={fontsize=\scriptsize}, listing only, colback=codebg, colframe=codebg, rounded corners}
// Открыть файл на диске
FILE *f1 = fopen("in.txt", ...);

// Стандартный ввод: то, что мы вводим с клавиатуры
FILE *f2 = stdin;

// Стандартный вывод: то, что печатается в терминале
FILE *f3 = stdout;
\end{tcblisting}
\end{frame}

\begin{frame}[fragile]
\frametitle{Текстовые и бинарные форматы}
\begin{itemize}
\item На физическом диске хранятся просто байты, их интерпретация -- вопрос конкретной программы и конкретного формата файла
\pause
\item Текстовые форматы:
\pause
\begin{itemize}
\item Интерпретируются, как последовательность текстовых символов: например, число \mintinline{cpp}|100| представляется тремя байтами \mintinline{cpp}|'1', '0', '0'|
\pause
\item Часто содержит спецсимволы: пробелы, табуляцию, перевод строк
\pause
\item {\color{green}+} Проще читать человеку
\pause
\item {\color{red}\texttt{-}} Сложнее читать программе, больше размер файла
\end{itemize}
\pause
\item Бинарные форматы:
\pause
\begin{itemize}
\item Строго следуют определённому порядку байт
\pause
\item Каждому байту задана его интерпретация, например число \mintinline{cpp}|100| может представляться как один байт \mintinline{cpp}|0x64|
\pause
\item {\color{green}+} Проще читать программе, компактный размер файла
\pause
\item {\color{red}\texttt{-}} Сложно читать человеку
\end{itemize}
\end{itemize}
\end{frame}

\begin{frame}[fragile]
\frametitle{Работа с файлами}
\usemintedstyle{lightbulb}
\begin{tcblisting}{listing engine=minted, minted language=cpp, minted options={fontsize=\scriptsize}, listing only, colback=codebg, colframe=codebg, rounded corners}
FILE* fopen(char *path, char *mode);
\end{tcblisting}
\pause
\usemintedstyle{tango}
\begin{itemize}
\item \mintinline{cpp}|path| -- путь до файла
\pause
\item \mintinline{cpp}|mode| -- режим открытия файла, может состоять из нескольких символов:
\pause
\begin{itemize}
\item Один из \mintinline{cpp}|'r'/'w'/'a'| -- открыть на чтение, на перезапись, или на дозапись в конец (при записи файл создаётся, если не существовал)
\pause
\item \mintinline{cpp}|'+'| -- режим ``обновления'' файла: можно и читать, и писать
\pause
\item \mintinline{cpp}|'x'| -- выдать ошибку, если файл открыт на запись (\mintinline{cpp}|'w'/'a'|) но не существует
\pause
\item \mintinline{cpp}|'b'| -- открыть файл в бинарном режиме (по умолчанию режим текстовый)
\pause
\item \mintinline{cpp}|'t'| -- \textbf{нестандартное расширение}: открыть файл в текстовом режиме
\end{itemize}
\end{itemize}
\end{frame}

\begin{frame}[fragile]
\frametitle{Текстовый и бинарный режимы}
\begin{itemize}
\item В бинарном режиме файл читается/пишется, как есть, без преобразований
\pause
\item В текстовом режиме под Windows:
\pause
\begin{itemize}
\item При чтении перевод строки \mintinline{cpp}|"\r\n"| превращается в \mintinline{cpp}|"\n"|, при записи -- наоборот
\pause
\item Особое значение байта \mintinline{cpp}|0x1a| означает конец файла
\end{itemize}
\pause
\item В текстовом режиме под другими платформами:
\pause
\begin{itemize}
\item Ничего особого не происходит :)
\end{itemize}
\pause
\item Не важно при работе с текстовыми форматами
\pause
\item При работе с бинарными форматами \textbf{обязательно} нужно указывать флаг \mintinline{cpp}|"b"|
\end{itemize}
\end{frame}

\begin{frame}[fragile]
\frametitle{Работа с файлами}
\usemintedstyle{lightbulb}
\begin{tcblisting}{listing engine=minted, minted language=cpp, minted options={fontsize=\scriptsize}, listing only, colback=codebg, colframe=codebg, rounded corners}
FILE *file = fopen("my_file.txt", "r");
if (file == NULL) {
  // Обработаем ошибку открытия файла
}

// Работаем с файлом...

// Закрываем файл
fclose(file);
\end{tcblisting}
\pause
\usemintedstyle{tango}
\begin{itemize}
\item Зачем закрывать файл, если при закрытии программы ОС его закроет сама?
\pause
\item На \mintinline{cpp}|struct FILE| тратится память
\pause
\item Число файловых дескрипторов на один процесс ограничено
\pause
\item Открытый файл нельзя удалить
\end{itemize}
\end{frame}

\begin{frame}[fragile]
\frametitle{Бинарные файлы}
\usemintedstyle{lightbulb}
\begin{tcblisting}{listing engine=minted, minted language=cpp, minted options={fontsize=\scriptsize}, listing only, colback=codebg, colframe=codebg, rounded corners}
// Прочитать count объектов, каждый по size байт,
// из файла в буфер dst_buf
size_t fread(void *dst_buf, size_t size, size_t count, FILE *file);

// Записать count объектов, каждый по size байт,
// из буфера src_buf в файл
size_t fwrite(void *src_buf, size_t size, size_t count, FILE *file);
\end{tcblisting}
\pause
\usemintedstyle{tango}
\begin{itemize}
\item Почему не просто количество байт?
\pause
\item Функции возвращают количество прочитанных/записанных \textit{объектов}
\end{itemize}
\end{frame}

\begin{frame}[fragile]
\frametitle{Бинарные файлы}
\usemintedstyle{lightbulb}
\begin{tcblisting}{listing engine=minted, minted language=cpp, minted options={fontsize=\scriptsize}, listing only, colback=codebg, colframe=codebg, rounded corners}
// Открыли на чтение в бинарном режиме
FILE *file1 = fopen("in.bin", "rb");

// Открыли на запись в бинарном режиме
FILE *file2 = fopen("out.bin", "wb");

int array[100];

// Прочитали 100 int'ов из файла в массив
fread(array, sizeof(int), 100, file1);

// Записали 100 int'ов из массива в файл
fwrite(array, sizeof(int), 100, file2);
\end{tcblisting}
\end{frame}

\begin{frame}[fragile]
\frametitle{Текстовые файлы}
\usemintedstyle{lightbulb}
\begin{tcblisting}{listing engine=minted, minted language=cpp, minted options={fontsize=\scriptsize}, listing only, colback=codebg, colframe=codebg, rounded corners}
// Открыли на чтение в текстовом режиме
FILE *file1 = fopen("in.text", "r");

// Открыли на запись в текстовом режиме
FILE *file2 = fopen("out.txt", "w");

int i = 0;
float f = 0.f;

// Прочитать целое число и число с плавающей
// точкой из файла в текстовом формате
fscanf(file1, "%i %f", &i, &f);

// Записать их в другой файл, тоже в 
// текстовом формате
fprintf(file2, "%i %f", i, f);
\end{tcblisting}
\pause
\usemintedstyle{tango}
\begin{itemize}
\item Тоже возвращают количество прочитанных/записанных \textit{объектов}
\end{itemize}
\end{frame}

\begin{frame}[fragile]
\frametitle{Текстовые файлы}
\usemintedstyle{lightbulb}
\begin{tcblisting}{listing engine=minted, minted language=cpp, minted options={fontsize=\scriptsize}, listing only, colback=codebg, colframe=codebg, rounded corners}
char string[256];
fscanf(file, "%s", str);
\end{tcblisting}
\begin{itemize}
\item Прочитает всё до первого \textit{whitespace} символа: пробела/табуляции/перевода строки
\pause
\item В чём проблема такого кода? \pause Мы не знаем заранее размер строки, и \mintinline{cpp}|fscanf| не знает размер нашего буфера
\pause
\item Альтернативы:
\pause
\usemintedstyle{tango}
\begin{itemize}
\item \mintinline{cpp}|fgets(string, 255, file)|
\pause
\item \mintinline{cpp}|fscanf(file, "%255s", string)|
\pause
\item \mintinline{cpp}|fscanf_s(file, "%s", string, 255)|
\end{itemize}
\end{itemize}
\end{frame}

\begin{frame}[fragile]
\frametitle{printf/fprintf/sprintf}
\begin{itemize}
\item Функции \mintinline{cpp}|printf/scanf| работают со \textit{стандартными потоками} ввода/вывода (stdin/stdout)
\pause
\item Функции \mintinline{cpp}|fprintf/fscanf| работают с произвольными файлами
\pause
\item Функции \mintinline{cpp}|sprintf/sscanf| работают со строками
\end{itemize}
\pause
\usemintedstyle{lightbulb}
\begin{tcblisting}{listing engine=minted, minted language=cpp, minted options={fontsize=\scriptsize}, listing only, colback=codebg, colframe=codebg, rounded corners}
int a, b;
char str1[] = "3 4";
sscanf(str1, "%d %d", &a, &b);

char str2[256];
sprintf(str2, "%d + %d = %d", a, b, c);
\end{tcblisting}
\end{frame}

\begin{frame}[fragile]
\frametitle{Конец файла, ошибки}
\usemintedstyle{lightbulb}
\begin{tcblisting}{listing engine=minted, minted language=cpp, minted options={fontsize=\scriptsize}, listing only, colback=codebg, colframe=codebg, rounded corners}
while (true) {
  fread(...);

  if (feof(file)) {
    // Конец файла
    break;
  }

  if (ferror(file)) {
    // Обработаем ошибку
  }

  ...
}
\end{tcblisting}
\pause
\usemintedstyle{tango}
\begin{itemize}
\item \mintinline{cpp}|feof(file)| возвращает состояние флага EOF (End Of File)
\pause
\item \mintinline{cpp}|ferror(file)| возвращает состояние ошибок (узнать саму ошибку можно через \mintinline{cpp}|errno|, про это поговорим позже)
\pause
\item Флаги выставляются \textbf{при попытке чтения}: сначала мы попробуем прочитать за концом файла, и только потом выставится флаг EOF
\end{itemize}
\end{frame}

\begin{frame}[fragile]
\frametitle{Буферизация, flush}
\usemintedstyle{lightbulb}
\begin{tcblisting}{listing engine=minted, minted language=cpp, minted options={fontsize=\scriptsize}, listing only, colback=codebg, colframe=codebg, rounded corners}
int main() {
  FILE *file = fopen("file.txt", "w");
  fprintf(file, "%i", 42);

  // Произошла какая-то ошибка
  if (some_error) {
    return 0;
  }

  // Нормальное завершение
  fclose(file);
  return 0;
}
\end{tcblisting}
\usemintedstyle{tango}
\begin{itemize}
\item В чём проблема кода? \pause Записанные данные лежат в промежуточном буфере в \mintinline{cpp}|struct FILE|, и могут не попасть на диск
\pause
\item Можно явно вызвать \mintinline{cpp}|fflush(file)| -- это форсирует сброс данных на диск
\end{itemize}
\end{frame}

\end{document}
